\documentclass[10pt]{amsart}
\usepackage[a4paper,margin=23mm]{geometry}
\usepackage{amsmath,amssymb,amsthm,mathtools,hyperref,graphicx,xspace,letltxmacro,xpatch}
\usepackage[english]{babel}
\usepackage{xfrac}

\usepackage[scr=rsfs,cal=euler]{mathalfa}
\multlinegap0pt

\newcommand{\Psfrac}[2]{{(\sfrac{#1}{#2})}}
\newcommand{\psfrac}[2]{{\sfrac{(#1)}{#2}}}
\newcommand{\spfrac}[2]{{\sfrac{#1}{(#2)}}}
\newcommand{\pspfrac}[2]{{\sfrac{(#1)}{(#2)}}}

\newcommand{\Babs}[1]{\Bigl\vert#1\Bigr\vert}
\newcommand{\Bnorm}[1]{\Bigl\Vert#1\Bigr\Vert}

\usepackage{csquotes}

\usepackage{mathtools,stmaryrd,esint}

\usepackage[shortlabels]{enumitem}
\usepackage[nameinlink]{cleveref}
\crefdefaultlabelformat{#2\textup{#1}#3}

\usepackage{longtable,tabu}
\usepackage{aligned-overset}

\theoremstyle{plain}
\newtheorem{theorem}{Theorem}[section]
\newtheorem{lemma}[theorem]{Lemma}
\newtheorem{proposition}[theorem]{Proposition}
\newtheorem{corollary}[theorem]{Corollary}

\theoremstyle{definition}
\newtheorem{definition}[theorem]{Definition}
\newtheorem{remark}[theorem]{Remark}
\newtheorem{example}[theorem]{Example}

\newcommand{\dd}{\mathop{}\mathopen{}\mathrm{d}}
\newcommand{\C}{\mathbb{C}}
\newcommand{\N}{\mathbb{N}}
\newcommand{\Q}{\mathbb{Q}}
\newcommand{\R}{\mathbb{R}}
\newcommand{\Z}{\mathbb{Z}}
\newcommand{\Sph}{\mathbb{S}}
\newcommand{\Cspace}{\mathscr{C}}
\newcommand{\Mspace}{\mathscr{M}}
\newcommand{\ener}{\mathcal{E}}
\newcommand{\mass}{\boldsymbol{M}}
\newcommand{\mres}{\mathop{\hbox{\vrule height 6pt width .5pt depth 0pt \vrule height .5pt width 4pt depth 0pt}}\nolimits}
\newcommand{\id}{\mathrm{Id}}
\newcommand{\loc}{\mathrm{loc}}
\newcommand{\lsc}{\mathrm{lsc}}
\newcommand{\eps}{\varepsilon}
\newcommand{\uloc}{L^1_\mathrm{uloc}}
\newcommand{\st}{\;:\;}

\DeclarePairedDelimiter{\abs}{\lvert}{\rvert}
\DeclarePairedDelimiter{\norm}{\lVert}{\rVert}
\DeclarePairedDelimiterX{\intint}[2]{\llbracket}{\rrbracket}{#1,#2}
\DeclareMathOperator{\supp}{supp}
\DeclareMathOperator{\diam}{diam}
\DeclareMathOperator{\tang}{Tan}
\DeclareMathOperator*{\esssup}{ess\,sup}
\DeclareMathOperator{\hdm}{\mathcal{H}}
\DeclareMathOperator{\lbm}{\mathcal{L}}
\DeclareMathOperator{\dist}{dist}
\DeclareMathOperator{\reach}{reach}

\newcommand{\weakto}{\rightharpoonup}
\newcommand{\mweakto}{\xrightharpoonup{\Cspace_0'}}
\newcommand{\narrowto}{\xrightharpoonup{\Cspace_b'}}
\newcommand{\xto}[2][]{\xrightarrow[#2]{#1}}
\newcommand{\xmweakto}[1]{\xrightharpoonup[#1]{\Cspace_0'}}
\newcommand{\xweakto}[1]{\xrightharpoonup[#1]{}}
\newcommand{\xnarrowto}[1]{\xrightharpoonup[#1]{\Cspace_b'}}


\usepackage{xfrac,xcolor,graphicx}\graphicspath{{./figures/}}\providecommand{\skpt}{\par\smallskip\noindent}\providecommand{\pointir}{.\ }\providecommand{\eg}{e.g.}\providecommand{\ie}{i.e.}\providecommand{\resp}{resp.}\providecommand{\cf}{cf.}\providecommand{\etc}{etc.}\let\Subsection\subsection\let\Subsubsection\subsubsection
\let\oldbigcup\bigcup
